\documentclass[twocolumn]{aastex631}
\begin{document}
\title{The reionization ionizing-photon-budget shortfall for star-forming galaxies is not robust to systematics at z\~{}6 (optimistic systematic corner)}
\author{NebulaMind Lab (autonomous pipeline)}
\affiliation{NebulaMind Lab, \url{https://lab.nebulamind.net}}
\begin{abstract}
Reionization ionizing-photon-budget reconciliation at z\~{}6: star-forming galaxies require f\_esc=0.013 (+0.011/-0.006) to close the budget (Madau-Dickinson SFRD, log xi\_ion=25.5+/-0.15, clumping C=2-5, JWST-SFRD tail), versus indirect-proxy-inferred f\_esc=0.080 (+0.147/-0.051) from LzLCS O32/beta calibrations. Median delta(required-inferred)=-0.065 dex-frac (16-84\%: -0.212 to -0.014); 6\% of the systematic MC shows a shortfall. Conclusion: the budget CLOSES within the systematic, and the sign holds under both O32 and beta calibrations. The result is bounded by the xi\_ion x clumping x proxy-calibration systematic, not by statistics. Generated autonomously from public data (jwst) via the NebulaMind Lab runner: a bounded, reproducible, descriptive study.
\end{abstract}
\section{Introduction}
The question of whether star-forming galaxies alone can provide sufficient ionizing photons to drive reionization remains a topic of ongoing debate [Muñoz2024]. Recent studies have highlighted potential discrepancies between the observed galaxy population and the required photon budget, suggesting that additional sources or revised assumptions may be necessary [Davies2021]. To address this issue, we revisit the reionization-photon-budget using literature-anchored values for key parameters.
\section{Data and method}
In our analysis, we adopt the cosmic star formation rate density (SFRD) from Madau \& Dickinson's (2014) analytic fitting function. We also use published calibrations for the ionizing photon production efficiency (xi\_ion) and the escape fraction (f\_esc) proxy based on O32/beta ratios [Chisholm+22, Flury+22; Simmonds+24]. Notably, we do not rely on new observational data from surveys like JWST or SDSS, but instead focus on reconciling existing literature values to assess the photon budget.
\section{Result}
Our calculation reveals that star-forming galaxies can close the reionization-photon-budget at z\~{}6 if they have an escape fraction of f\_esc=0.013 (+0.011/-0.006). This value is lower than the indirect-proxy-inferred escape fraction of 0.080 (+0.147/-0.051) derived from LzLCS O32/beta calibrations. The median difference between the required and inferred escape fractions is -0.065 dex-frac, with a range of -0.212 to -0.014 (16-84\% confidence interval). Notably, 6\% of our systematic Monte Carlo simulations indicate a shortfall in the photon budget.
\begin{figure}\centering\includegraphics[width=\columnwidth]{result.png}\caption{The reionization ionizing-photon-budget shortfall for star-forming galaxies is not robust to systematics at z\~{}6 (optimistic systematic corner)}\end{figure}
\section{Caveats}
While our analysis provides a reconciliation of the reionization-photon-budget within the systematic uncertainties, it is essential to acknowledge the limitations of our approach. Our calculation relies on a single selection of literature values for key parameters, which may not fully capture the complexity and variability of the underlying astrophysical processes. Additionally, the use of uncalibrated proxy calibrations introduces uncertainty, as these relationships may not hold universally across all galaxy populations. Furthermore, our study does not account for potential contributions from other sources, such as active galactic nuclei or X-ray binaries, which could also influence the photon budget. A more comprehensive understanding will require further observational and theoretical efforts to refine these assumptions and constraints.
\section*{References}
\footnotesize
[Muoz2024] Muñoz et al. (2024). Reionization after JWST: a photon budget crisis?. bibcode 2024MNRAS.535L..37M \par [Davies2021] Davies et al. (2021). The Predicament of Absorption-dominated Reionization: Increased Demands on Ionizing Sources. bibcode 2021ApJ...918L..35D \par [Park2022] Park et al. (2022). Calibrating excursion set reionization models to approximately conserve ionizing photons. bibcode 2022MNRAS.517..192P \par [Duncan2015] Duncan et al. (2015). Powering reionization: assessing the galaxy ionizing photon budget at z < 10. bibcode 2015MNRAS.451.2030D \par [Madau2017] Madau (2017). Cosmic Reionization after Planck and before JWST: An Analytic Approach. bibcode 2017ApJ...851...50M

\end{document}
